Napjaink felgyorsult világában kevesen használják a tömegközlekedést, mert az nem elég kényelmes, nem elég gyors vagy éppen nem úgy működik, ahogy azt elvárnánk tőle. A gyakori problémák közé tartozik, hogy a buszok nem tartják a menetrendet vagy a jegyek megvásárlása nehézkes. Sok helyen a menetrend nincs kifüggesztve a megállókban és az interneten sem elérhető. Az is gyakran megesik hogy a buszok nagyon zsúfoltak és erről az utazni vágyók csak akkor vesznek tudomást, amikor már a megállóban vannak és nem áll rendelkezésre más utazási lehetőség.

A fent említett problémákra kínál megoldást a Bus4U online platform, amely egy mobilalkalmazásból, egy felhasználói weboldalból és egy adminisztrátori weboldalból tevődik össze. A végfelhasználók számára elérhető funkciók közé tartozik a buszjáratok keresése a kiindulópont, a célállomás és az időpont megadásával, a járatokra szóló előzetes online jegyvásárlás, illetve a járatok indulási időpontjainak megtekintése városokra és megállókra lebontva. Ezen kívül megtekinthetőek a buszmegállók a térképen, városok és útvonalak szerint szűrve, valamint valós időben követhetőek az autóbuszok pillanatnyi helyzetei is.

Az alkalmazás másik fontos része az adminisztrációs felület, amelyet a busztársaságok használhatnak. Ezen a felületen lehetőség van a menetrendek kezelésére, új útvonalak és megállók felvitelére, valamint a meglévő adatok frissítésére. Továbbá itt kezelhetők az alkalmazottak információi, beleértve a sofőrök és egyéb személyzet adatait, valamint az autóbuszok műszaki és üzemeltetési adatai. Ez az adminisztrációs rendszer biztosítja a busztársaságok adatainak karbantartását, így azok könnyedén frissíthetik szolgáltatásaikat és biztosíthatják a pontosságot.

A rendszer különböző szolgáltatásainak biztosításához, mint például a jegykezelés, szükség van egy másik alkalmazásra is, amely a buszokon található POS terminálon fut. Ez biztosítja a jegyek érvényesítését, illetve az autóbuszok élő információinak megosztását is a központi szerverrel.

Továbbfejlesztési lehetőségek között szerepelnek a bérletek és különböző kedvezmények létrehozása, valamint a buszok telítettségének követése valós időben. A busztársaságok szempontjából fontos lehet még a sofőrök munkaidejének követése és az automatizált könyvelés bevezetése is.


\textbf{Kulcsszavak:} tömegközlekedés, buszjárat, digitalizálás